% !TEX root = ../TCC - CD e IA.tex
\chapter{Referencial Teórico}\label{Referencial Teórico}

\section{Mídias Sociais e Marketing Político Digital}

A transição do marketing político para a era digital representa uma mudança de paradigma fundamental, abandonando o modelo tradicional, unidirecional e focado no "produto" \cite{KAID2004}, para adotar uma lógica interativa e dialógica \cite{FREITAS2021}. Essa evolução, que teve seu marco inicial no processo de redemocratização brasileiro e foi catalisada globalmente por campanhas pioneiras como a de Barack Obama \cite{NEWMAN2016}, consolidou a internet como o campo de batalha central pela opinião pública \cite{ITUASSU2023}. A comunicação política, antes mediada por instituições tradicionais \cite{COOK2005}, agora ocorre em um ecossistema de redes onde a interação direta e a construção de comunidades se tornam objetivos primordiais, reconfigurando a relação entre governantes e governados \cite{LIPSCHULTZ2022}.

O impacto do uso destas tecnologias nos resultados eleitorais é inegável, como demonstram as campanhas "hipermidiáticas" de Donald Trump e Jair Bolsonaro, que se caracterizaram pelo uso estratégico de big data e pela disseminação de desinformação \cite{ITUASSU2023}. A vitória de Bolsonaro em 2018, em particular, quebrou o paradigma da comunicação política no Brasil, superando a necessidade de tempo de TV e grande estrutura partidária com uma campanha digital-first \cite{IASULAITIS2022, FERNANDES2018}. O WhatsApp emergiu como a ferramenta decisiva nesse pleito, sendo o principal vetor para a viralização de conteúdo e mobilização de apoiadores \cite{ITUASSU2023}. A influência é tão concreta que uma pesquisa do DataSenado (2019) revelou que 45\% dos eleitores brasileiros decidiram seu voto para presidente em 2018 com base em informações vistas nas redes sociais.

Contudo, a ascensão das mídias sociais também intensificou desafios complexos para a democracia  \cite{NETO2024}. A arquitetura das plataformas fomenta a criação de "bolhas" e "câmaras de eco"  \cite{BARBERA2020}, que, por sua vez, alimentam uma polarização predominantemente afetiva, baseada na hostilidade ao campo adversário \cite{SILVEIRA2020}. Esse ambiente polarizado torna-se um terreno fértil para a disseminação de desinformação \cite{GUESS2020}, fenômeno que se tornou evidente durante a pandemia de Covid-19, quando a politização da crise sanitária e a circulação de conteúdos falsos andaram lado a lado \cite{GOULART2022, RECUERO2021}. O combate a esse fenômeno tornou-se um dos principais desafios para a integridade do processo democrático \cite{TRASEL2024}.

\section{Marketing de Crescimento e Análise de Engajamento Digital}

Se a seção anterior estabeleceu o ambiente no qual a comunicação política digital se realiza, cabe agora examinar como esse ambiente é medido e otimizado. O campo do \emph{marketing} de crescimento (\emph{growth marketing}) oferece um arcabouço consolidado para essa tarefa, definido como a ciência e a arte de extrair, gerir e analisar grandes volumes de dados digitais \textemdash{} textuais, imagéticos e de rede \textemdash{} para converter o comportamento do público em decisões acionáveis \cite{FEROZ2024}. Trata-se de um deslocamento de ênfase: da mera mensuração da atividade para a orientação sistemática da ação, na qual cada métrica coletada deve estar vinculada a um objetivo de negócio ou de comunicação.

Um conceito central nesse arcabouço é a distinção entre métricas de vaidade (\emph{vanity metrics}) e métricas acionáveis. As primeiras \textemdash{} como o número absoluto de curtidas, seguidores ou visualizações \textemdash{} são de fácil obtenção e produzem uma impressão de sucesso, mas raramente informam uma decisão, pois não distinguem, por exemplo, o engajamento motivado pela aclamação daquele movido pela controvérsia. As segundas vinculam-se diretamente a resultados: taxas de conversão, retenção e crescimento efetivo da audiência. Essa distinção fundamenta a proposição de uma \emph{North Star Metric} (métrica-norte), um indicador único que sintetiza o valor entregue ao público e alinha toda a estratégia em torno de um objetivo prioritário, em contraposição ao acompanhamento disperso de dezenas de indicadores superficiais \cite{FEROZ2024}.

Para organizar as etapas pelas quais um público é conquistado e retido, a literatura de \emph{marketing} digital consolidou modelos de funil. O modelo AIDA \textemdash{} atenção, interesse, desejo e ação \textemdash{} descreve a progressão clássica da persuasão, ao passo que o modelo RACE \textemdash{} \emph{Reach} (alcance), \emph{Act} (ação), \emph{Convert} (conversão) e \emph{Engage} (engajamento) \textemdash{} adapta essa lógica ao ambiente digital, oferecendo estágios mensuráveis que conectam a descoberta de um conteúdo à construção de um relacionamento duradouro \cite{FEROZ2024}. Cada estágio do funil admite indicadores próprios: o alcance é aferido por impressões e visualizações; a ação, por curtidas e cliques; a conversão, por comportamentos que sinalizam adesão; e o engajamento, pela recorrência e pela produção espontânea de conteúdo pelo próprio público.

A transposição desses conceitos \textemdash{} originalmente formulados para o \emph{marketing} comercial \textemdash{} ao domínio político é direta em sua estrutura, ainda que exija reinterpretação de seus termos. No contexto de um perfil político, o alcance corresponde à disseminação de uma publicação; a ação, às interações de baixo custo, como curtidas; a conversão, à manifestação de apoio explícito nos comentários; e o engajamento, à mobilização de apoiadores que amplificam espontaneamente a mensagem. Essa reinterpretação fornece o esqueleto sobre o qual, mais adiante, este trabalho articulará os resultados de sua análise, convertendo métricas de engajamento e sinais de sentimento em um diagnóstico orientado à decisão de assessorias políticas.

\section{Engajamento do Público e Atividades Online Relacionadas à Marca (COBRAs)}

Os modelos de funil descritos na seção anterior organizam os estágios do engajamento, mas não explicam suas motivações \textemdash{} isto é, por que um indivíduo decide interagir com um conteúdo. Para essa dimensão, a literatura de comportamento do consumidor consolidou o construto das \emph{Consumers' Online Brand-Related Activities} (COBRAs), proposto por Muntinga, Moorman e Smit \cite{MUNTINGA2011}. É digno de nota que uma das autoras seminais do conceito atua no campo da comunicação política, o que evidencia a proximidade original entre o construto e o domínio deste trabalho e legitima sua transposição.

O modelo COBRAs organiza o engajamento com conteúdo em três níveis crescentes de intensidade. O primeiro, o \emph{consumo}, corresponde à participação passiva: assistir a um vídeo, visualizar uma imagem ou ler uma publicação sem produzir qualquer conteúdo. O segundo, a \emph{contribuição}, envolve a interação com o conteúdo existente e com outros usuários, por meio de curtidas, compartilhamentos e, sobretudo, comentários. O terceiro e mais intenso, a \emph{criação}, ocorre quando o indivíduo produz e publica conteúdo original relacionado à marca \cite{MUNTINGA2011}. Essa gradação é diretamente análoga aos estágios de um funil de engajamento e fornece a base conceitual para interpretar os diferentes tipos de interação observados em um perfil político: o comentário, por exemplo, situa-se no nível de contribuição, um patamar de engajamento superior ao mero consumo.

Investigações posteriores buscaram identificar os valores que motivam a passagem de um nível a outro. No contexto de campanhas virais, Dinh e Lee \cite{DINHLEE2024} propuseram um modelo de seis valores \textemdash{} facilidade de navegação, valor hedônico, valor funcional, valor estético, interação social e autoidentidade \textemdash{} que motivam o público a consumir, contribuir e criar. Tais valores oferecem um vocabulário útil para explicar por que certos conteúdos falham em engajar (por exemplo, por deficiência de valor hedônico ou estético), enquanto outros mobilizam intensamente o público. Complementarmente, a dimensão afetiva do engajamento tem sido estudada por meio do construto do \emph{brand love}: pesquisas demonstram que o vínculo emocional com uma marca é um antecedente do engajamento, mediando a disposição do público para consumir, comentar e criar conteúdo, efeito por sua vez moderado pelo valor percebido da marca \cite{CASTROGONZALEZ2025}. Transposta ao domínio político, essa dimensão afetiva encontra paralelo na lealdade e no vínculo emocional entre o eleitorado e a figura pública, ajudando a explicar por que perfis distintos suscitam padrões de engajamento tão heterogêneos.

A aplicação desses construtos ao ambiente político deve, contudo, ser feita com cautela metodológica. Os modelos de COBRAs e de \emph{brand love} foram originalmente formulados para marcas comerciais, e sua transposição para figuras políticas exige reconhecer as diferenças de contexto \textemdash{} notadamente o fato de que o engajamento político é frequentemente movido por polarização afetiva, na qual a hostilidade ao campo adversário é um motor tão potente quanto a adesão. Reconhecida essa ressalva, o construto permanece valioso: ele fornece a este trabalho a lente teórica para interpretar os comentários não como meras ocorrências de sentimento, mas como manifestações de engajamento em diferentes níveis, articuláveis aos estágios do funil de crescimento.

\section{Precedentes de Análise de Audiência Política em Escala}

A articulação entre \emph{marketing} de crescimento, construtos de engajamento e comunicação política, proposta neste trabalho, não é inédita: encontra precedente consolidado na literatura de \emph{big social data}. O estudo de Flesch, Vatrapu e Mukkamala \cite{FLESCH2017} sobre a eleição federal alemã de 2017 constitui a referência metodológica mais próxima desta pesquisa. Os autores analisaram a interação de aproximadamente um milhão de indivíduos com as páginas dos sete principais partidos alemães, demonstrando a viabilidade de aplicar, em escala, técnicas de segmentação de audiência e métricas de crescimento a dados de engajamento político.

Três aspectos desse precedente são particularmente relevantes para o presente trabalho. O primeiro é a segmentação da audiência em categorias funcionais \textemdash{} leais (\emph{loyalists}), defensores (\emph{advocates}), críticos (\emph{critics}) e ativistas (\emph{activists}) \textemdash{} derivadas dos padrões de interação, e não de atributos declarados. Essa segmentação demonstra que os próprios comportamentos de engajamento permitem identificar tipos de público estrategicamente distintos, um princípio que fundamenta o uso da clusterização como ferramenta de segmentação nesta pesquisa. O segundo é o emprego de métricas de crescimento tipicamente associadas ao \emph{marketing}, como a taxa de crescimento mensal composta (\emph{Compound Monthly Growth Rate}, CMGR) e a análise de retenção da audiência ao longo do tempo, aplicadas de forma legítima ao domínio político. O terceiro é a lógica comparativa: ao contrastar os sete partidos entre si, os autores mostram que o desempenho digital só adquire significado quando avaliado de forma relativa \textemdash{} princípio que, neste trabalho, justifica o tratamento dos vinte e sete governadores como parâmetros recíprocos de comparação (\emph{benchmarking}).

A abordagem de Flesch e colaboradores insere-se, ademais, em uma tradição metodológica mais ampla de análise de mídias sociais para a comunicação política, que estrutura o processo analítico em etapas sucessivas de preparação, análise, rastreamento e descoberta. Essa estruturação processual dialoga diretamente com a arquitetura de dados adotada neste trabalho, na qual a coleta, o tratamento e a modelagem são organizados em camadas sucessivas e reprodutíveis. Em conjunto, esses precedentes conferem respaldo acadêmico à proposta central desta pesquisa: a de que sinais de engajamento e sentimento extraídos de redes sociais podem ser organizados, por meio de vocabulário e métricas de crescimento, em um instrumento de diagnóstico e decisão para a comunicação política.

\section{Fundamentos de Processamento de Linguagem Natural (PLN)}

O campo do Processamento de Linguagem Natural (PLN) foi fundamentado por contribuições teóricas que estabeleceram os alicerces para a interação entre computadores e a linguagem humana. Alan Turing, embora não fosse um linguista, propôs um conceito fundamental com seu "Teste de Turing", que desafiou a máquina a exibir um comportamento inteligente indistinguível do de um ser humano, estabelecendo um objetivo visionário para a área \cite{TURING1950}. Pouco depois, o linguista Noam Chomsky revolucionou a linguística com a teoria da gramática gerativa, apresentada em sua obra "Syntactic Structures", que postulava um conjunto finito de regras capaz de gerar todas as sentenças gramaticalmente corretas de um idioma \cite{CHOMSKY1957}. Essa abordagem influenciou profundamente a primeira era do PLN, que se baseava em sistemas simbólicos e regras gramaticais codificadas manualmente. Um exemplo prático e icônico desse período foi o programa ELIZA, desenvolvido por Joseph Weizenbaum. ELIZA simulava um psicoterapeuta rogeriano através de simples scripts de reconhecimento de padrões e substituição de palavras-chave, demonstrando pela primeira vez a capacidade de um programa de manter uma conversa aparentemente coerente e evidenciando o quão prontamente os humanos atribuem inteligência às máquinas \cite{WEIZENBAUM1966}.

A partir do final da década de 1980, as limitações dos sistemas baseados em regras para lidar com a ambiguidade e a complexidade da linguagem real impulsionaram uma mudança de paradigma. A chamada "revolução estatística" emergiu, trocando as regras gramaticais por modelos probabilísticos derivados de grandes volumes de texto (corpora). Pesquisadores da IBM, como Frederick Jelinek e Peter F. Brown, foram pioneiros nessa abordagem, especialmente no campo da tradução automática, publicando trabalhos seminais que demonstravam a superioridade dos métodos estatísticos \cite{BROWN1990}. Técnicas como os Modelos Ocultos de Markov (HMMs), popularizadas em tutoriais influentes como o de Lawrence Rabiner, tornaram-se padrão para tarefas como a etiquetagem morfossintática e o reconhecimento de fala \cite{RABINER1989}. A consolidação desse conhecimento ocorreu com a publicação de livros didáticos que se tornaram referência, como "Foundations of Statistical Natural Language Processing", de Christopher Manning e Hinrich Schütze, que educaram uma geração de pesquisadores nos novos métodos quantitativos \cite{MANNING1999}.

A era contemporânea do PLN é dominada pelo deep learning (aprendizagem profunda), que começou a ganhar força na década de 2010. Uma das inovações cruciais foi o desenvolvimento de representações de palavras em vetores densos, os word embeddings, com destaque para o modelo Word2Vec de Tomas Mikolov e sua equipe no Google, que capturava relações semânticas entre palavras de forma eficiente \cite{MIKOLOV2013}. Arquiteturas de redes neurais como as LSTMs (Long Short-Term Memory), propostas originalmente por Sepp Hochreiter e Jürgen Schmidhuber, permitiram que os modelos lidassem melhor com dependências de longo prazo em sequências de texto \cite{HOCHREITER1997}. No entanto, a maior disrupção veio com a publicação do artigo "Attention Is All You Need", que introduziu a arquitetura Transformer. Essa nova arquitetura, baseada exclusivamente em mecanismos de atenção, permitiu o processamento paralelo de texto e uma captura de contexto mais eficaz, superando as limitações das redes recorrentes \cite{VASWANI2017}. O Transformer tornou-se a base para os Grandes Modelos de Linguagem (LLMs), inaugurando uma nova fase de avanços com modelos como o BERT, que introduziu o pré-treinamento bidirecional \cite{DEVLIN2018}.

Atualmente, o campo avança em um ritmo acelerado, impulsionado por equipes de pesquisa em grandes empresas de tecnologia e na academia. Os autores dos principais LLMs, como o GPT-3 da OpenAI \cite{BROWN2020} e o LLaMA da Meta AI \cite{TOUVRON2023}, continuam a expandir as fronteiras do que é possível em geração e compreensão de linguagem. Figuras como Dan Jurafsky e Christopher Manning permanecem extremamente influentes, atualizando constantemente suas obras de referência, como "Speech and Language Processing", para incorporar os últimos avanços. No Brasil, o campo de PLN também possui uma comunidade ativa e crescente. Iniciativas como o grupo "Brasileiras em PLN" visam aumentar a visibilidade e a colaboração entre pesquisadoras na área. Centros de pesquisa consolidados, como o Centro de Inteligência Artificial (C4AI) e o Núcleo Interinstitucional de Linguística Computacional (NILC) da USP, são polos de excelência que contribuem com pesquisa de ponta e para a formação de novos talentos, garantindo a participação do Brasil no cenário global de inovação em Processamento de Linguagem Natural.

\section{Técnicas de Mineração de Texto para Análise de Mídias Sociais}

A mineração de texto é definida como uma extensão da mineração de dados, focada na extração de informações úteis e desconhecidas de documentos textuais em linguagem natural. \cite{PEZZINI2023} No contexto das mídias sociais, essa técnica assume uma complexidade adicional devido à natureza não estruturada, informal e ambígua dos dados, onde as pessoas frequentemente não se preocupam com a ortografia ou com a construção gramatical correta das sentenças. \cite{IRFAN2015} A literatura de pesquisa tem explorado a aplicação de técnicas de mineração de texto para compreender o comportamento de interação pessoa-para-pessoa, identificar padrões de pensamento humano e agrupar usuários. \cite{IRFAN2015} Pesquisadores brasileiros como Anderson Pezzini, Christian Aranha e Aydano P. Machado têm contribuído para esta área, aplicando a mineração de texto em contextos como análise de artigos científicos e educação a distância. \cite{PEZZINI2023} A pesquisa aponta para uma carência na literatura de estudos que se concentrem especificamente em conjuntos de dados de redes sociais, o que confere ao projeto do TCC, com seu foco no Instagram, uma relevância particular para preencher essa lacuna.

\subsection{Análise de Sentimentos}

A Análise de Sentimento, também conhecida como Mineração de Opinião, foi consolidada como um campo de pesquisa a partir do trabalho seminal de Bo Pang e Lillian Lee, publicado em 2008. \cite{PANG2008} Eles definiram a tarefa como o processo de identificar a subjetividade ou a polaridade de um texto, classificando a linguagem com base em sua carga emocional (negativa, positiva ou neutra). A contribuição de Pang e Lee foi fundamental não apenas por fornecer uma definição clara, mas por explorar as implicações mais amplas da análise de opinião, como questões de privacidade e o potencial de manipulação.

A literatura acadêmica distingue as duas abordagens principais para a análise de sentimento: a baseada em léxicos e a baseada em aprendizado de máquina (ML) e grandes modelos de linguagem (LLMs). \cite{VANDERVEENBLEICH2025} A abordagem baseada em léxicos utiliza dicionários de palavras pré-rotuladas com pontuações de sentimento. Suas vantagens incluem generalização e independência de domínio, permitindo a análise de textos em diferentes contextos sem a necessidade de treinamento prévio. No entanto, pode sacrificar algum desempenho em comparação com métodos mais complexos. Já uma Abordagem Baseada em ML/LLMs utiliza modelos treinados em grandes volumes de dados para classificar o sentimento. Esses modelos podem alcançar uma acurácia superior dentro de domínios específicos, mas exigem mais recursos computacionais e podem apresentar limitações de generalização e viés.\cite{VANDERVEENBLEICH2025} A escolha do TCC de utilizar o modelo cardiffnlp/twitter-xlm-roberta-base-sentiment reflete a decisão de priorizar a alta acurácia e a capacidade de capturar as nuances semânticas dos textos de mídias sociais, uma vez que se trata de uma abordagem moderna baseada em transformers.

Os modelos transformer representam uma mudança de paradigma na área de PLN, notadamente pela introdução do mecanismo de self-attention. Diferentemente de arquiteturas anteriores, como as Redes Neurais Recorrentes (RNNs) e as Long Short-Term Memory networks (LSTMs), que processavam sequências de texto de forma linear (uma palavra por vez), os transformers permitem o processamento paralelo. Isso significa que eles podem considerar o contexto de uma palavra em relação a todas as outras palavras na sequência de uma só vez, capturando dependências sintáticas e semânticas de maneira mais eficaz. \cite{GOWDA2025} Essa arquitetura aprimorada se manifesta em modelos como o BERT (Bidirectional Encoder Representations from Transformers) e, mais recentemente, o RoBERTa (Robustly Optimized BERT Approach). O RoBERTa, desenvolvido pelo Facebook AI, é uma versão otimizada do BERT que aprimora o processo de pré-treinamento. Ele foi treinado em um conjunto de dados 10 vezes maior que o do BERT (~160 GB vs. ~16 GB), com um vocabulário estendido (50 mil tokens vs. 30 mil tokens) e sem a tarefa de "previsão da próxima sentença" (NSP), que foi considerada ineficaz. O RoBERTa também utiliza uma estratégia de mascaramento dinâmico, onde as palavras a serem mascaradas mudam a cada época de treinamento. Essas otimizações resultaram em um desempenho superior do RoBERTa em diversas tarefas de PLN, incluindo análise de sentimentos. \cite{PRASANTHI2023} A escolha do TCC pelo modelo cardiffnlp/twitter-xlm-roberta-base-sentiment é estrategicamente relevante, pois ele é uma variação multilíngue do XLM-RoBERTa, pré-treinado em cerca de 198 milhões de tweets. Por ter sido treinado especificamente em dados de mídias sociais, esse modelo é particularmente adequado para lidar com a linguagem informal e as nuances encontradas em comentários de plataformas como o Instagram, assegurando alta acurácia na classificação de sentimentos.

\subsection{Modelagem de Tópicos}

O LDA (Latent Dirichlet Allocation), proposto por David Blei, Andrew Ng e Michael Jordan em 2003, é a técnica de modelagem de tópicos mais influente na literatura acadêmica. É um modelo probabilístico generativo que assume que os documentos são representados como misturas aleatórias sobre tópicos latentes, e que cada tópico é caracterizado por uma distribuição sobre as palavras. Embora revolucionário, o LDA opera sob a suposição de "saco de palavras" (que a ordem das palavras em um documento pode ser negligenciada) e não captura a semântica contextual, o que o torna menos eficaz para analisar textos curtos e informais, como os de mídias sociais. \cite{BLEINGJORDAN2003}

O artigo "BERTopic: Neural topic modeling with a class-based TF-IDF procedure", do autor Maarten Grootendorst, apresenta uma nova abordagem para a modelagem de tópicos, chamada BERTopic, para que se possa superar as limitações dos modelos de tópicos tradicionais, como o LDA, ao integrar técnicas de clustering e uma variação do TF-IDF baseada em classes para gerar representações de tópicos coerentes. \cite{GROOTENDORST2022}

O BERTopic apresenta várias vantagens em relação ao LDA, destacando-se principalmente pela sua flexibilidade e capacidade de gerar tópicos coerentes. Uma das principais inovações do BERTopic é a separação entre o processo de embedding de documentos e a representação de tópicos. Isso permite que diferentes procedimentos de pré-processamento sejam aplicados em cada etapa, o que não é tão facilmente alcançado com o LDA. Essa separação confere ao modelo uma maior adaptabilidade, possibilitando o uso de modelos de linguagem de última geração para criar embeddings, enquanto ainda permite a utilização de modelos clássicos em situações onde recursos computacionais são limitados. \cite{GROOTENDORST2022}

Além disso, o BERTopic utiliza uma abordagem baseada em clustering, que se mostra mais eficaz em capturar a semântica dos documentos. Ao empregar técnicas como UMAP para redução de dimensionalidade antes do agrupamento com HDBSCAN, o BERTopic consegue formar clusters de documentos que representam tópicos distintos de maneira mais precisa. Essa abordagem resulta em uma maior coerência dos tópicos gerados, além de permitir a modelagem dinâmica de tópicos, algo que o LDA, sendo um modelo estático, não suporta adequadamente \cite{GROOTENDORST2022}. Portanto, o BERTopic não apenas melhora a qualidade dos tópicos identificados, mas também oferece uma maior flexibilidade e adaptabilidade em sua aplicação.

\subsection{Clusterização}

Formalmente, a análise de cluster pode ser definida como a tarefa de particionar um conjunto de objetos em grupos, de modo que os objetos dentro de um mesmo grupo sejam mais similares entre si do que com aqueles em outros grupos \cite{JAINMURTYFLYNN1999, STATHOPOULOU2019}. Cada objeto é tipicamente representado como um vetor de medições ou um ponto em um espaço multidimensional. O objetivo primário, portanto, é maximizar a homogeneidade intra-cluster (similaridade interna) e, simultaneamente, maximizar a heterogeneidade inter-cluster (dissimilaridade externa) \cite{TANSTEINBACHKUMAR2006}. Uma característica distintiva e central da clusterização é sua natureza de aprendizado não supervisionado. Diferentemente da classificação supervisionada, na análise de cluster, a afiliação dos objetos aos grupos é desconhecida a priori; não existem rótulos pré-definidos para guiar o processo de aprendizado \cite{XUWUNSCH2005}. O algoritmo deve, portanto, descobrir essas estruturas "naturais" ou "ocultas" nos dados de forma autônoma, tornando-se uma ferramenta indispensável para a análise exploratória de dados \cite{GKOLLORISTUDI2023}. Historicamente, as raízes da análise de cluster são profundas e interdisciplinares. Suas origens remontam à antropologia, com os trabalhos de Driver e Kroeber em 1932, sendo posteriormente introduzida na psicologia por Zubin em 1938 e Tryon em 1939 \cite{DRIVERKROEBER1932, ZUBIN1938, TRYON1939}. Essa trajetória demonstra a necessidade perene, em diversas áreas do conhecimento, de organizar observações em agrupamentos coerentes para facilitar a compreensão e a tomada de decisão.

A aplicação da análise de cluster pode ser categorizada em dois papéis principais, embora frequentemente sobrepostos: clusterização para compreensão e clusterização para utilidade \cite{TANSTEINBACHKUMAR2006} A distinção entre esses dois propósitos é fundamental, pois influencia diretamente a escolha do algoritmo, dos parâmetros e das métricas de avaliação.
Clusterização para Compreensão visa descobrir classes ou grupos conceitualmente significativos que revelem a estrutura natural inerente aos dados. Neste contexto, os clusters são vistos como representações de categorias reais no domínio do problema. Clusterização para Utilidade, por outro lado, emprega o agrupamento como um passo intermediário para um objetivo prático, onde os clusters em si não precisam corresponder a classes naturais, mas devem ser úteis para a tarefa final.

A vasta literatura sobre análise de cluster produziu centenas de algoritmos, cada um com diferentes noções do que constitui um cluster e como encontrá-los eficientemente. Esses algoritmos podem ser organizados em uma taxonomia baseada em seus princípios operacionais subjacentes. As principais famílias de algoritmos incluem métodos particionais, hierárquicos, baseados em densidade e baseados em modelo. Os métodos particionais visam dividir o conjunto de dados em um número pré-especificado de k clusters não sobrepostos, onde cada ponto de dados pertence a exatamente um grupo \cite{FAHAD2014}. A partição é tipicamente obtida através da otimização de uma função objetivo, geralmente baseada em uma medida de distância. Um exemplo de algoritmo que pertence a esta família é o K-means. Diferentemente dos métodos particionais, os algoritmos hierárquicos não produzem uma única partição dos dados. Em vez disso, eles constroem uma sequência aninhada de partições que pode ser visualizada como uma estrutura em árvore, conhecida como dendrograma  
\cite{HENRY2005}. O dendrograma ilustra como os clusters são fundidos (ou divididos) em diferentes níveis de similaridade, oferecendo uma visão rica e multi-resolução da estrutura dos dados.

\subsection{Abordagens Híbridas e Síntese do Referencial}

As três técnicas descritas nesta seção \textemdash{} análise de sentimentos, modelagem de tópicos e clusterização \textemdash{} são, com frequência, empregadas de forma isolada na literatura. Cada uma, tomada individualmente, responde a uma pergunta parcial: a análise de sentimentos revela \emph{como} o público reage, mas não sobre o quê; a modelagem de tópicos revela \emph{sobre o que} se discute, mas não com que valência; a clusterização agrupa por semelhança, mas nada diz sobre o conteúdo do discurso. A premissa central deste trabalho é que a integração sinérgica dessas três abordagens produz uma leitura da percepção pública superior à soma de suas partes, revelando padrões que nenhuma delas, isoladamente, é capaz de expor.

Essa integração, contudo, permaneceria descritiva não fosse sua articulação aos arcabouços de engajamento e crescimento discutidos no início deste capítulo. É nesse ponto que as camadas teóricas se encontram. Os resultados da clusterização, ao segmentarem publicações e perfis por padrões de desempenho, fornecem os grupos análogos aos segmentos de audiência identificados na literatura de \emph{big social data} \cite{FLESCH2017}. A análise de sentimentos qualifica cada um desses grupos, distinguindo o engajamento movido por aprovação daquele movido por controvérsia. A modelagem de tópicos, por fim, permite contrastar o discurso produzido pelas assessorias \textemdash{} nas legendas e transcrições \textemdash{} com os temas a que o público de fato reage, nos comentários.

O cruzamento dessas evidências pode, então, ser organizado segundo a estrutura de um funil de engajamento que articula os níveis dos COBRAs \cite{MUNTINGA2011} aos estágios do modelo RACE \cite{FEROZ2024}. Nessa articulação, o alcance de uma publicação corresponde ao estágio de descoberta; as curtidas e o consumo, ao engajamento passivo; os comentários, ao nível de contribuição; e a amplificação espontânea, ao nível de criação e engajamento mais profundo. É precisamente esse arcabouço integrado \textemdash{} que converte descrições analíticas em um diagnóstico orientado à decisão \textemdash{} que fundamenta a metodologia e a aplicação desenvolvidas nos capítulos seguintes, nos quais a proposta é operacionalizada e validada no estudo de caso dos governadores brasileiros.


